\documentclass[10pt]{article}
\usepackage[T1]{fontenc}
\usepackage[utf8]{inputenc}
\usepackage[french]{babel}
\usepackage{newpxtext,newpxmath}
\usepackage[a4paper,top=1.5cm,bottom=1.35cm,left=1.6cm,right=1.6cm]{geometry}
\usepackage{microtype}
\usepackage{xcolor}
\usepackage{titlesec}
\usepackage[most]{tcolorbox}
\usepackage{enumitem}
\usepackage{array}
\usepackage{booktabs}
\usepackage{fancyhdr}

\definecolor{navy}{HTML}{1F3A5F}
\definecolor{grey}{HTML}{54606E}
\definecolor{green}{HTML}{2C6E49}
\definecolor{amber}{HTML}{9C6A16}
\definecolor{red}{HTML}{A6322E}
\definecolor{plum}{HTML}{7A4EA8}
\definecolor{boxgray}{HTML}{F3F5F8}
\definecolor{boxred}{HTML}{FCF2F1}
\definecolor{boxgreen}{HTML}{F1F7F3}
\definecolor{boxplum}{HTML}{F7F3FB}
\definecolor{boxcode}{HTML}{F7F8FA}
\definecolor{boxedge}{HTML}{D9E0EA}

\setlength{\parindent}{0pt}
\setlength{\parskip}{2.6pt}
\linespread{1.0}

\pagestyle{fancy}\fancyhf{}
\lhead{\small\color{grey}Scella}
\rhead{\small\color{grey}Les couples -- étude détaillée -- 11 août 2026}
\cfoot{\small\color{grey}\thepage}
\renewcommand{\headrulewidth}{0.3pt}

\titleformat{\section}{\color{navy}\large\bfseries}{}{0pt}{}
\titlespacing{\section}{0pt}{11pt}{3pt}
\titleformat{\subsection}{\color{navy}\bfseries}{}{0pt}{}
\titlespacing{\subsection}{0pt}{7pt}{1.5pt}

\newtcolorbox{these}{colback=boxgray,colframe=boxedge,boxrule=0.5pt,arc=2.5pt,
  left=8pt,right=8pt,top=4.5pt,bottom=4.5pt}
\newtcolorbox{alerte}{colback=boxred,colframe=red!55,boxrule=0.8pt,arc=2.5pt,
  left=8pt,right=8pt,top=4.5pt,bottom=4.5pt}
\newtcolorbox{acquis}{colback=boxgreen,colframe=green!45,boxrule=0.8pt,arc=2.5pt,
  left=8pt,right=8pt,top=4.5pt,bottom=4.5pt}
\newtcolorbox{piste}{colback=boxplum,colframe=plum!45,boxrule=0.8pt,arc=2.5pt,
  left=8pt,right=8pt,top=4.5pt,bottom=4.5pt}
\newtcolorbox{extrait}{colback=boxcode,colframe=boxedge,boxrule=0.5pt,arc=1.5pt,
  left=7pt,right=7pt,top=3.5pt,bottom=3.5pt}

\newcommand{\best}[1]{{\color{green}\textbf{#1}}}
\newcommand{\pire}[1]{{\color{red}\textbf{#1}}}
\newcommand{\sz}[1]{{\small\color{grey}\textsf{#1}}}
\newcommand{\code}[1]{\texttt{\small #1}}
\newcommand{\x}{$\times$}
\newcommand{\loc}[1]{{\footnotesize\color{grey}\texttt{#1}}}

% En-tête normalisé de chaque couple
\newcommand{\couple}[4]{%
  \subsection{#1 · #2 \hfill {\normalsize\color{grey}moyenne \textbf{#3} · écart #4}}}

\begin{document}

\begin{center}
{\Huge\bfseries\color{navy} Scella}\\[4pt]
{\large\color{grey}\itshape Les huit couples : anatomie, mécanisme, falsifieur, notation}\\[4pt]
{\small\color{grey} 11 août 2026 \; · \; étude détaillée \; · \; branche \code{couche2-m0-frontier}, HEAD \code{580782f}}
\end{center}
\vspace{2pt}

\begin{these}
\textbf{Ce que ce document ajoute aux trois précédents.} Les premiers disaient \emph{quels} couples
existent, \emph{dans quel ordre} ils valent, et \emph{contre qui} ils tiennent. Celui-ci descend d'un
cran : pour chacun, le \textbf{mécanisme exact} -- souvent une ligne -- le \textbf{test qui échoue}
quand la propriété casse, et la \textbf{justification de chacune des quatre notes} plutôt que la note
seule. Deux précisions gagnées en descendant : l'injectivité du couple n\textsuperscript{o}2 tient
entièrement dans un \code{Option::take()}, et la taxonomie de dégradation du couple
n\textsuperscript{o}7 compte \textbf{dix motifs nommés} -- une rigueur qui rend d'autant plus
frappante l'absence de tout producteur.
\end{these}

\section{1 · Méthode de notation}

Un couple est deux éléments \textbf{plus un liant} : le mécanisme qui les fait produire ensemble une
propriété qu'aucun n'a seul. Chaque couple est noté sur les quatre lentilles intrinsèques établies le
5 août -- sans lentille d'objectif extérieur, celle-ci ayant disparu avec la piste qu'elle servait.

\textbf{Maturité} : ce qu'il reste à faire pour que la propriété soit réellement obtenue, caveats
compris. \textbf{Rareté} : la paire existe-t-elle ailleurs, ou la discipline qu'elle exige est-elle
inhabituelle. \textbf{Adversaire} : la propriété survit-elle à un examinateur hostile qui cherche la
faille plutôt que la confirmation. \textbf{Démontrabilité} : peut-on la montrer à un ingénieur
sceptique, et surtout peut-on lui montrer son \emph{falsifieur} échouer.

Les rangs vont de 1 (meilleur) à 8, par lentille. La moyenne des quatre donne le classement ;
l'\textbf{écart} entre le meilleur et le pire rang mesure l'instabilité, et il se lit séparément.

\begin{alerte}
\textbf{Faits et jugements, séparés.} Les mécanismes, emplacements, noms de tests et taxonomies
ci-dessous sont des \textbf{faits}, lus dans le code du 11 août et cités avec leur ligne. Les
\textbf{notes} sont un \textbf{jugement} structuré pour être discutable. Un désaccord sur « rareté = 3 »
est légitime ; un désaccord sur « \code{take\_matching} consomme le record » se tranche en ouvrant le
fichier.
\end{alerte}

\renewcommand{\arraystretch}{1.12}
\begin{center}\small
\begin{tabular}{@{}r >{\raggedright\arraybackslash}p{6.6cm} cccc c c@{}}
\toprule
& \textbf{Le couple} & \sz{Mat} & \sz{Rar} & \sz{Adv} & \sz{Dém} & \textbf{Moy} & \sz{Éc.} \\
\midrule
\textbf{1} & Artefact scellé \x{} re-prouveur zéro-dépendance & \best{1} & 3 & \best{1} & 3 & \textbf{2,00} & 2 \\
\textbf{2} & Topologie déclarée \x{} appariement injectif & 2 & 4 & 2 & 2 & \textbf{2,50} & 2 \\
\textbf{3} & Liaison v4 \x{} ancre publiée & 3 & 5 & 3 & \best{1} & \textbf{3,00} & 4 \\
\textbf{4} & Canari de Lipschitz \x{} porte de démarrage & 6 & 2 & 4 & 4 & \textbf{4,00} & 4 \\
5\,= & Caviardage PII \x{} scellement & 4 & 7 & 6 & 6 & 5,75 & 3 \\
5\,= & Plafonnement \x{} coût scellé & 5 & \pire{8} & 5 & 5 & 5,75 & 3 \\
7 & Co-signature R3 \x{} chaîne HMAC & \pire{8} & \best{1} & \pire{8} & \pire{8} & 6,25 & \pire{7} \\
8 & Revendication \x{} sa propre chute & 7 & 6 & 7 & 7 & 6,75 & \best{1} \\
\bottomrule
\end{tabular}
\end{center}

\section{2 · Les huit couples, en détail}

\couple{Rang 1}{Artefact scellé \x{} re-prouveur zéro-dépendance}{2,00}{2}

\textbf{Anatomie.} \emph{A} : l'artefact scellé -- \code{meta\_steps.jsonl} et
\code{meta\_manifest.json}, les octets exacts produits par le chemin servi. \emph{B} : un re-prouveur
dont l'arbre de dépendances normales est \textbf{vide}. \emph{Le liant} : l'absence de chaîne
d'approvisionnement entre l'auditeur et l'octet.

\textbf{Le mécanisme.} Atteindre zéro dépendance interdit \code{serde}, \code{sha2}, \code{hmac} --
tout. Les deux crates ré-implémentent donc à la main SHA-256, HMAC et un parseur JSON :
\code{verify-mini} en \textbf{1\,159 lignes} sur sept modules
\sz{(\code{sha256}, \code{hmac}, \code{json}, \code{record}, \code{verify}, \code{lib}, \code{main})},
\code{verify-meta-mini} en \textbf{992 lignes} sur six. Ce n'est pas un accident de graphe de
dépendances : c'est une contrainte tenue \textbf{deux fois}, sur deux piles indépendantes.

Ré-implémenter de la cryptographie est une mauvaise idée par défaut, et le dépôt le sait : la
\emph{fidélité} des deux implémentations est prouvée séparément, par des \textbf{vecteurs golden gelés}
et un \textbf{cross-check octet-à-octet} contre la référence -- \code{golden\_chain\_hmac\_matches\_}
\code{audit\_core}, \code{roundtrip\_against\_efe\_meta\_verify}, plus un aller-retour par version de
schéma \sz{(v2 \code{safety\_delta}, v3 \code{vote\_hashes})}. La ré-implémentation est donc encadrée
par une preuve d'équivalence, pas laissée à la confiance.

\textbf{Le falsifieur.} \code{cargo tree -p verify-meta-mini -e normal -{}-prefix none} doit rendre
\emph{une seule ligne}. C'est le contrôle le plus court du dépôt et le plus difficile à contester.

\textbf{La notation.} \best{Maturité 1} : le seul couple du portefeuille dont l'analyse n'a relevé
\emph{aucun} caveat. \sz{Rareté 3} : la discipline est rare, mais le \emph{principe} d'un vérifieur
minimal ne l'est pas. \best{Adversaire 1} : la revendication se vérifie en une commande, et aucune
rhétorique n'y résiste. \sz{Démontrabilité 3} : excellente, mais moins spectaculaire que le vert
gênant du rang 3 -- une ligne de sortie convainc un ingénieur, elle ne le surprend pas.

\begin{alerte}
\textbf{La faille de ce couple n'est pas dans son code, elle est dans sa garde.} L'invariant
« \code{cargo tree} ne montre que ce crate » est \emph{documenté} en commentaire dans quatre
\code{Cargo.toml}, mais n'est gardé par \textbf{aucun test et aucun job d'intégration continue}. Le
seul contrôle exécutable vit dans \code{scripts/demo-consensus-receipt.sh}, que rien ne déclenche
automatiquement. Un \code{cargo add} passerait la revue sans qu'aucun voyant ne s'allume, et le couple
le mieux classé du portefeuille tomberait en silence. \sz{Correctif : porter le contrôle en test.
Taille XS.}
\end{alerte}

\couple{Rang 2}{Topologie déclarée \x{} appariement injectif}{2,50}{2}

\textbf{Anatomie.} \emph{A} : la topologie \textbf{déclarée} -- la liste des appels annoncés avant
exécution. \emph{B} : un \code{LlmProvenanceRecord} scellé par appel \textbf{réellement} effectué.
\emph{Le liant} : un appariement \textbf{injectif} entre les deux.

\textbf{Le mécanisme, et il tient en une ligne.} L'appariement puise dans un vivier de records qui
sont des \code{Option}, et \emph{consomme} celui qu'il retient \loc{(mission\_graph.rs:98)} :

\begin{extrait}\small
\code{fn take\_matching(pool: \&mut [Option<LlmProvenanceRecord>], provider: Option<\&str>)}\\
\hspace*{2em}\code{...}\\
\hspace*{2em}\code{pool[idx].take()} \hfill \sz{$\leftarrow$ toute l'injectivité est ici}
\end{extrait}

\code{Option::take()} laisse un \code{None} derrière lui : un record ne peut donc être apparié
\textbf{qu'une seule fois}. Sans cette ligne, trois n\oe uds déclarés sur le même fournisseur
pointeraient tous le \emph{même} enregistrement et le reçu conclurait « cohérent » -- un faux positif
parfait, et indétectable, puisque toutes les empreintes seraient valides.

Ce n'est pas un cas théorique. La topologie \code{llm/debate} déclare $2 \cdot \text{rounds} + 1$
appels avec les mêmes fournisseurs \textbf{répétés} à chaque ronde \loc{(mission\_graph.rs:273)} ;
c'est exactement là que l'appariement naïf se casserait, et le dépôt l'a épinglé. Cinq topologies sont
servies, chacune avec son arité déclarée : \code{vote} \sz{($k$ = nombre de modèles)}, \code{relay}
\sz{(3 : planificateur, implémenteur, synthétiseur)}, \code{debate} \sz{($2r+1$)}, \code{specialist}
\sz{(nombre de sous-tâches)}, \code{cascade}.

\textbf{Les falsifieurs, et il y en a deux.}

\begin{extrait}\small
\code{omitted\_call\_makes\_verify\_false} \loc{(:475)} --- déclare trois fournisseurs, n'en exécute
que deux ; exige \code{fully\_proven == false} et \code{bridge.broken == 1}.\\[2pt]
\code{injective\_matching\_prevents\_partial\_receipt\_forgery} \loc{(:492)} --- déclare trois appels
au \emph{même} fournisseur, n'en produit qu'un ; exige un n\oe ud apparié et deux manquants.
\end{extrait}

Le second est celui qui compte : il couvre le cas qu'une implémentation raisonnable raterait.

\textbf{Le verdict est décomposé}, et c'est une qualité rare. Le premier test assère aussi
\code{v.swarm.ok()} -- la chaîne d'essaim reste \emph{intègre} pendant que \code{fully\_proven} est
faux. Le reçu ne dit donc pas seulement \emph{que} quelque chose ne va pas, il dit \emph{où} : ici le
pont vers les appels externes, pas le sceau. Un auditeur qui reçoit un verdict binaire doit enquêter ;
celui-ci reçoit un diagnostic.

\textbf{La notation.} \sz{Maturité 2} : servi, sans caveat propre -- mais bordé par le faux couple F2
\sz{(le \code{config\_hash} déclaré et non rapporté)}. \sz{Rareté 4} : déclarer une topologie est
banal, l'apparier injectivement l'est beaucoup moins. \sz{Adversaire 2} : l'attaque évidente est
précisément celle qui est couverte. \sz{Démontrabilité 2} : deux tests nommés, lisibles, et dont
l'échec est parlant.

\textbf{La limite, à énoncer.} Ce couple prouve la \emph{conformité à la déclaration}, pas
l'\emph{ambition} de la déclaration. Un manifeste qui annonce honnêtement un panel d'un seul modèle
produit un reçu parfaitement valide.

\couple{Rang 3}{Liaison v4 \x{} ancre publiée}{3,00}{4}

\textbf{Anatomie.} \emph{A} : la \code{RunBinding} v4 -- identifiants de participants \textbf{ordonnés},
\code{config\_hash} déclarés alignés, opérateur, $\alpha$, $\gamma$, scellés ensemble. \emph{B} : un
journal d'ancres append-only portant la racine \textbf{et le nombre de pas}. \emph{Le liant} :
l'\textbf{antériorité} de la publication sur la réponse.

\textbf{Le mécanisme, côté v4.} L'encodage canonique est binaire et préfixé en longueur
\loc{(efe-meta-verify/src/lib.rs:165)}, avec deux protections que seule une lecture attentive révèle.
D'abord, le préfixe de comptage n'est \emph{pas} la longueur de la première colonne :

\begin{extrait}\small
\code{let n = self.participant\_ids.len().min(self.declared\_config\_hashes.len());}
\end{extrait}

Une désérialisation contournant le constructeur pourrait produire des colonnes désalignées ; le
préfixe mentirait alors sur le nombre d'entrées qui suivent, et l'encodage deviendrait ambigu. Le
\code{min} ferme les deux bouts. Ensuite le \textbf{marqueur de présence} du \code{config\_hash}
est scellé lui aussi \sz{(un octet 1 ou 0)} : sans lui, « absent » et « présent avec un hash nul » se
confondraient.

\textbf{Le mécanisme, côté ancre.} La ronde publie sa racine \emph{avant} que la réponse ne soit
construite \loc{(controller/iterate.rs:280 $\rightarrow$ :291)}, et un échec de publication fait
échouer l'appel. La clé du journal est préfixée \code{"meta:"} \loc{(anchor.rs:81)}, de sorte qu'un
\textbf{seul} journal couvre la Couche 2 et la Couche 3, et qu'un auditeur sache à quoi il a affaire
en lisant la clé.

\textbf{Pourquoi l'ancre est nécessaire, et non décorative.} Le préfixe d'une chaîne HMAC valide est
lui-même une chaîne HMAC valide. Retirer les derniers pas et « corriger » le manifeste produit donc un
bundle qui se re-prouve \textbf{parfaitement} : rien, à l'intérieur, ne trahit la coupure. Seule une
racine publiée à l'avance l'expose -- et le contrôle ne demande \textbf{aucune clé}.

Le contrôle compare le \textbf{nombre de pas d'abord}, et l'ordre est délibéré
\loc{(verify-meta-mini/src/verify.rs:369)} : une troncature change les deux valeurs, mais « le journal
engage 30 pas, le bundle en présente 29 » est un diagnostic, là où « les racines diffèrent » ne dit à
l'auditeur que ce qu'il voyait déjà. Le message distingue même \code{TRONCATURE} d'\code{AJOUT}.

\textbf{Le falsifieur -- et c'est le meilleur du dépôt, parce qu'il est \emph{gênant}.} La
démonstration montre d'abord le bundle tronqué \textbf{passer} la re-preuve avec la clé -- un vert
franc, mérité, et trompeur -- avant de le faire tomber sur l'ancre sans aucune clé. Montrer sa propre
garantie échouer là où elle ne s'applique pas est ce qui rend crédible l'endroit où elle s'applique.

\textbf{La notation.} \sz{Maturité 3} : servi et complet, mais l'ancre reste détenue par l'opérateur.
\sz{Rareté 5} : l'ancrage est un principe connu \sz{(journaux de transparence)} ; c'est le
\emph{couplage} avec une liaison nominative qui est moins courant. \sz{Adversaire 3} : très solide
contre la troncature rétroactive, attaquable sur la détention du journal. \best{Démontrabilité 1} : le
meilleur du portefeuille, pour la raison ci-dessus.

\textbf{La limite.} L'ancre est chez l'opérateur. Elle supprime la réécriture \emph{rétroactive}, pas
la confiance. La garantie tierce est le jalon 3b, et \code{verify\_anchor} existe déjà côté
\code{forensic-export}.

\couple{Rang 4}{Canari de Lipschitz \x{} porte de démarrage}{4,00}{4}

\textbf{Anatomie.} \emph{A} : un estimateur de constante de Lipschitz \textbf{empirique}, en boîte
noire \loc{(participants/lipschitz.rs)}. \emph{B} : le participant efe-core réel et sa porte de
démarrage \loc{(participant/local.rs:342)}. \emph{Le liant} : le \textbf{fail-closed} -- la mesure ne
produit pas un rapport, elle produit un refus de servir.

\textbf{Le mécanisme.} L'estimateur \sz{(\code{estimate\_lipschitz})} échantillonne des paires de
points et cherche adversarialement le pire rapport, par bissection, afin de traquer les
discontinuités. Il rend un certificat :

\begin{extrait}\small
\code{NonExpansiveCertificate::Empirical \{ l\_hat, margin, n\_pairs\_scored, inconclusive \}}\\[2pt]
\sz{accepté seulement si} \code{l\_hat < 1 − margin} \sz{\textbf{et} concluant.}
\end{extrait}

\textbf{Ce que la documentation dit d'elle-même, et qu'il faut citer.} \emph{« $\hat{L}$ est une borne
\textbf{inférieure} sur $L_\Phi$ »} ; le certificat est \emph{property-tested}, \textbf{pas} certifié ;
et un passage est \emph{« a falsification failure, never a proof »}. Le dispositif ne prouve donc
jamais la non-expansivité -- il échoue à la réfuter. C'est une distinction que peu de systèmes
prennent la peine d'écrire, et elle est ici dans le type lui-même via \code{inconclusive}.

\textbf{L'actif réel n'est pas le garde-fou, c'est ce que la mesure a produit.} Elle a
\textbf{réfuté le certificat analytique} du document de conception : $\hat{L} = 0{,}69$ en topologie
dynamique et $0{,}366$ en topologie fixe, contre un $\beta_1 \approx 0{,}17$ annoncé. La borne honnête
est la borne mesurée. Un dépôt qui publie la réfutation de sa propre note théorique dit quelque chose
sur sa méthode qu'aucune démonstration ne dirait.

\textbf{Un piège déjà tombé, et corrigé.} Une première version avalait les explosions de $\Phi$ comme
« hors domaine » via un \code{.ok()} -- un \emph{fail-open} parfait, où l'échec le plus grave
ressemblait au cas le plus bénin. Il a fallu un \code{PhiOutcome} à trois états pour distinguer
« hors domaine » de « a explosé ».

\textbf{La notation.} \sz{Maturité 6} : servi, mais avec une échappatoire déclarée. \sz{Rareté 2} :
mesurer empiriquement sa propre constante de contraction et refuser de démarrer dessus est très
inhabituel. \sz{Adversaire 4} : honnête sur son statut, mais l'échappatoire est une prise.
\sz{Démontrabilité 4} : le chiffre mesuré qui contredit le document est parlant, la démonstration est
longue \sz{(le sondage coûte de trente secondes à quelques minutes au démarrage)}.

\textbf{La limite.} \code{h1\_canary.mode = "off"} dans le manifeste fait servir le participant avec
un avertissement bruyant. L'échappatoire est déclarée, tracée, et elle existe.

\couple{Rang 5\,=}{Caviardage PII \x{} scellement}{5,75}{3}

\textbf{Anatomie.} \emph{A} : le caviardage selon la policy tenant, à quatre modes -- \code{Off},
\code{Detect}, \code{Redact}, \code{Strict} -- par tenant, avec \code{Off} pour défaut
\loc{(pii.rs:7)}. \emph{B} : le hash du prompt scellé dans la chaîne. \emph{Le liant} :
l'\textbf{ordre}.

\textbf{Le mécanisme.} \code{Off} et \code{Detect} ne caviardent pas ; \code{Redact} et \code{Strict}
caviardent, par détection \textbf{locale} -- expressions régulières, sans réseau. Le point décisif est
l'ordre d'appel : \code{redact\_for\_egress} s'exécute \textbf{avant tout usage} de \code{user\_text},
donc avant le hash scellé et avant le calcul d'arité. Le reçu commet ainsi ce qui est \emph{réellement}
sorti. Caviarder après scellement ne vaudrait rien : on scellerait le texte en clair et l'on jurerait
ensuite l'avoir protégé. Le patron servi est explicite dans le code :
\code{role $\rightarrow$ gate $\rightarrow$ clamp/redact $\rightarrow$ admit $\rightarrow$ scope
$\rightarrow$ run $\rightarrow$ seal}.

\textbf{La faille fermée, et elle était invisible.} Sans la feature \code{pii}, la ligne de caviardage
\emph{disparaissait du binaire} -- sans erreur, sans avertissement, sans trace d'exécution. Un
opérateur pouvait compiler, lire « préflight frontière OK », et expédier chaque prompt en clair. Le
correctif ne s'est pas contenté d'un \code{\#[cfg]} au site d'appel : la porte de démarrage exige
désormais \textbf{soit la capacité compilée, soit une reconnaissance explicite} de l'opérateur. Le
défaut \code{Off} par \emph{policy} reste légitime ; c'est le \textbf{silence} qui ne l'était pas.

\textbf{La notation.} \sz{Maturité 4} : servi, avec une fissure documentée. \pire{Rareté 7} :
caviarder avant un appel tiers est une pratique répandue ; seule la discipline d'ordre est
distinctive. \sz{Adversaire 6} : l'examinateur hostile trouve prise immédiatement. \sz{Démontrabilité
6} : montrer que l'\emph{ordre} compte exige de lire le code, ce qui est le contraire d'une
démonstration.

\textbf{La limite.} Sans policy tenant, le prompt sort en clair. C'est documenté dans le module même,
et c'est ce qui coûte à ce couple les trois rangs qu'il a perdus au reclassement.

\couple{Rang 5\,=}{Plafonnement \x{} coût scellé}{5,75}{3}

\textbf{Anatomie.} \emph{A} : \code{clamp\_budget} -- cap du tenant \textbf{puis} plafond serveur
\loc{(routes/llm.rs:249)}. \emph{B} : le coût par appel, scellé dans chaque enregistrement. \emph{Le
liant} : la borne posée \textbf{en amont} du sceau.

\textbf{Le mécanisme, et ses trois défaillances superposées corrigées.} La version antérieure ne
bornait rien en l'absence de politique tenant : \code{if let Some(pol) = ...} laissait passer intact
le budget d'un tenant sans policy. Aujourd'hui le plafond serveur s'applique \textbf{inconditionnellement},
après le cap tenant. Sa valeur vient de \code{EFE\_LLM\_MAX\_BUDGET\_CENTS}, par défaut
\textbf{10\,000}, avec un \code{.max(1)} explicite -- \emph{un plafond à zéro rendrait toute mission
impossible ; on refuse de se tirer dessus}. Et la conversion sature vers le budget \textbf{d'origine},
jamais vers \code{u32::MAX} : un dépassement arithmétique ne peut pas transformer un plafonnement en
autorisation.

\textbf{La notation.} \sz{Maturité 5} : servi et complet. \pire{Rareté 8} : le dernier du portefeuille
-- la comptabilité de coût dans une piste d'audit est un acquis de l'industrie. \sz{Adversaire 5} :
solide sur ce qu'il couvre. \sz{Démontrabilité 5} : convaincant, sans surprise.

\textbf{La limite.} Le couple ne protège pas contre un opérateur qui trafique sa propre table de prix
-- il scelle une dépense déclarée, pas une dépense auditée auprès du fournisseur. Le champ
\code{provider\_quoted\_cost\_cents} existe précisément pour ce recoupement.

\couple{Rang 7}{Co-signature R3 \x{} chaîne HMAC}{6,25}{7 \sz{(le plus instable)}}

\textbf{Anatomie.} \emph{A} : la chaîne HMAC tenant, qui scelle le record \textbf{co-signature
incluse} dans le \code{leaf\_hash}. \emph{B} : la signature Ed25519 du n\oe ud sur
\code{prompt\_hash} $\Vert$ \code{output\_hash}, clé privée ne quittant jamais le n\oe ud. \emph{Le
liant} : l'\textbf{orthogonalité} des deux systèmes.

\textbf{Le mécanisme, et l'argument qui le justifie est écrit dans le module.} Une co-signature HMAC
exigerait que le vérifieur \emph{détienne le secret du n\oe ud}, donc qu'il puisse le \textbf{forger} :
pas de non-répudiation. L'asymétrie est l'essence de R3. L'orchestrateur ne détient que du
\textbf{public} -- un tiers muni du registre, des records et de la clé tenant peut donc \textbf{tout
vérifier hors ligne, sans aucun secret}. D'où la propriété composée : \textbf{forger un R3 exige les
DEUX clés}.

\textbf{La règle de fer -- « aucune élévation sans preuve » (I-RC9).} Le barreau \textbf{scellé} est la
\emph{revendication} du producteur ; le barreau \textbf{vérifié} est ce que la preuve cryptographique
\emph{soutient}. Le vérifieur n'accepte jamais R3 sur parole : une co-signature invérifiable
\textbf{redescend} le record à R4 -- la transaction se re-prouve toujours -- et le verdict le
\textbf{signale}. C'est câblé, pas déclaratif : \code{n17\_chain.rs:346} calcule
\code{(held\_strength, cosign\_anomalies)}, la ligne 365 les place dans le verdict, la ligne 393
déclenche la dégradation, et \code{routes/llm.rs:201-202} sert \code{weakest\_rank} et
\code{weakest\_label} jusqu'au client.

\textbf{La taxonomie de dégradation compte dix motifs nommés}, et leur lecture donne la mesure du
soin : \code{NodeNotPinned}, \code{NodeRevoked}, \code{UnsupportedScheme}, \code{TupleMismatch},
\code{MalformedKey}, \code{MalformedSignature}, \code{SignatureInvalid}, \code{MissingOutput}
\sz{(refuser de lier une co-signature à un contenu absent)}, \code{NonceReused} \sz{(rejeu intra-run)}
et \code{ChainNotIntact} \sz{(pas de témoin sur un run non re-prouvé)}. Le registre distingue en outre
\code{Active} de \code{Revoked \{ reason \}}.

\begin{alerte}
\textbf{Et la propriété ne se déclenche jamais.} Aucun chemin servi ne \emph{produit} de co-signature :
les seules constructions de \code{NodeCoSignature} hors définitions de types sont des \code{\#[test]}
\loc{(efe-swarm-sdk/src/n17.rs:322)}, et le registre de n\oe uds n'est peuplé que dans des tests
\loc{(NodeRegistry::empty().with\_node(...), attest.rs:219)}. En production le registre est vide, aucun
enregistrement n'est R3, et le barreau le plus faible est \textbf{R4 par construction}. Dix motifs de
dégradation extrêmement soignés gardent une porte que personne ne franchit.
\end{alerte}

\textbf{La notation, et elle est brutale.} \pire{Maturité 8} : il manque un producteur.
\best{Rareté 1} : la conception -- double clé, non-élévation automatique, taxonomie à dix motifs -- est
la plus sophistiquée du portefeuille. \pire{Adversaire 8} : un examinateur le tue en une phrase, « donc
ça ne se déclenche jamais ». \pire{Démontrabilité 8} : il n'y a littéralement rien à montrer.

\begin{piste}
\textbf{L'écart de \pire{7} est le plus grand du tableau, et c'est l'information utile.} Un couple
premier sur une lentille et dernier sur trois n'est pas « moyen » : il est \textbf{instable}. La
formule juste est \emph{pire réalisé, meilleur latent}. Un producteur sur le chemin servi et un
registre peuplé le feraient passer de septième à disputer la première place -- c'est le seul chantier
identifié qui déplace un couple de six rangs. Et il a un précédent : le trou « aucun producteur » de la
pile B avait exactement cette forme, et il a été fermé.
\end{piste}

\couple{Rang 8}{Revendication \x{} sa propre chute}{6,75}{1 \sz{(le plus stable)}}

\textbf{Anatomie.} \emph{A} : les hops de reformulation, chacun portant son moteur effectif
\sz{(\code{SLM}, \code{RULE}, ou dégradé)}. \emph{B} : le prédicat « tous les hops sont SLM »
\loc{(routes/llm.rs:166)}. \emph{Le liant} : le calcul \textbf{depuis la sortie scellée}.

\textbf{Le mécanisme.} \code{slm\_mediated} vaut \code{!hops.is\_empty() \&\& slm == hops.len()} : la
revendication « relais médié par le petit modèle » n'est vraie que si \textbf{tous} les hops le sont,
et une seule dégradation la fait tomber, visiblement. La vue agrégée
\sz{(\code{hops}, \code{slm}, \code{rule}, \code{degraded}, \code{slm\_mediated})} est dérivée de la
sortie \textbf{scellée}, jamais recalculée ailleurs -- elle ne peut donc pas diverger de l'attesté.

\textbf{Une précision gagnée en descendant dans le code}, et qui renforce le couple : chaque hop est
\emph{en outre} scellé \textbf{individuellement} dans la chaîne HMAC, comme un n\oe ud Mode-A
\code{reform\_i} dont le JSON -- moteur effectif et faits d'audit -- est committé dans son
\code{output\_hash}, en plus du commitment global du blob de sortie par le n\oe ud \code{reduce}. Les
hops viennent de la sortie \textbf{productrice}, pas de la requête.

\textbf{La notation.} \sz{Maturité 7} : servi, mais voisin immédiat d'une fissure. \sz{Rareté 6} : le
patron « la revendication embarque son falsifieur » est élégant et peu répandu, mais sa portée est
étroite. \sz{Adversaire 7} : l'examinateur passe au faux couple F1 dans la phrase suivante.
\sz{Démontrabilité 7} : un booléen dans une réponse JSON.

\textbf{L'écart de 1 en fait le couple le plus \emph{stable} du portefeuille} -- il est régulièrement
avant-dernier partout. C'est l'inverse exact du rang 7, et les deux méritent d'être lus ensemble : la
stabilité n'est pas une qualité en soi, elle dit seulement qu'aucun travail ne le déplacera beaucoup.

\textbf{La limite, et c'est un faux couple voisin.} \code{slm\_mediated} atteste la \emph{médiation},
pas la \emph{souveraineté}. L'URL du petit modèle n'est pas contrainte au local, son egress n'est ni
gardé par l'allowlist tenant ni décompté \sz{(le hop est budgété 0 ¢)}, et le système avertit au
démarrage plutôt que de prétendre une souveraineté qu'il ne vérifie pas. Pire, la règle de
souveraineté existe en \textbf{trois} implémentations qui ont déjà divergé une fois en contournement
SSRF.

\section{3 · L'élément porteur -- ce que la décomposition révèle}

Décomposer huit couples en seize éléments fait apparaître qu'\textbf{ils ne sont pas seize}. Plusieurs
sont le même, réutilisé -- et l'un porte presque tout.

\renewcommand{\arraystretch}{1.15}
\begin{center}\small
\begin{tabular}{@{}>{\raggedright\arraybackslash}p{4.9cm} c >{\raggedright\arraybackslash}p{9.7cm}@{}}
\toprule
\textbf{L'élément} & \textbf{Couples} & \textbf{Ce que sa perte emporterait} \\
\midrule
\textbf{La chaîne HMAC scellée} & \pire{7 / 8} &
  1, 2, 3, les deux 5=, 7 et 8 \textbf{simultanément} : le point porteur unique \\
La policy tenant & 2 & les deux couples 5= -- caviardage \emph{et} plafond de coût \\
La sortie scellée, source unique & 2 & le couple 8, et l'appariement du couple 2 \\
La signature Ed25519 & 1 \sz{(+1 latent)} & le couple 7, et le chantier de signature de racine \\
\midrule
\textbf{L'estimateur de Lipschitz} & \best{1} & le couple 4 \textbf{seul} \\
\bottomrule
\end{tabular}
\end{center}

Le constat se vérifie en une commande : \code{participant/local.rs}, qui porte le canari, ne contient
\textbf{aucune} occurrence de HMAC. Le couple 4 est le seul qui survivrait intégralement à une faille
dans le scellement -- parce qu'il ne prouve rien à personne : il refuse de démarrer.

\begin{alerte}
\textbf{La conséquence réordonne les priorités.} Sept couples sur huit reposent sur le \emph{même}
élément. Une faille dans le scellement ne dégraderait pas le portefeuille, elle l'\textbf{effondrerait
d'un coup} -- y compris les mieux classés, y compris ceux dont la propriété semble indépendante. Ce
n'est pas un défaut de conception : c'est le prix d'une architecture cohérente où tout se ramène à un
sceau unique. Mais cela doit être \emph{su}.
\end{alerte}

Et cela donne au rang 1 la justification qui lui manquait. Si un seul élément porte sept couples, la
seule chose qui compte vraiment est de pouvoir \textbf{contrôler cet élément sans dépendre de celui qui
l'a produit}. Le re-prouveur zéro-dépendance n'est donc pas un « amplificateur » : c'est le
\textbf{seul contrôle indépendant du point porteur unique}. Ce qui rend d'autant plus gênant qu'aucun
test ne garde sa contrainte.

\section{4 · Sensibilité -- ce qui déplacerait chaque rang}

\renewcommand{\arraystretch}{1.2}
\begin{center}\small
\begin{tabular}{@{}r >{\raggedright\arraybackslash}p{4.3cm} >{\raggedright\arraybackslash}p{7.4cm} >{\raggedright\arraybackslash}p{3.4cm}@{}}
\toprule
& \textbf{Le couple} & \textbf{L'événement unique qui le déplacerait} & \textbf{Effet} \\
\midrule
1 & Re-prouveur zéro-dép & \emph{aucun} qui le monte -- il est au plafond ; une dépendance ajoutée le
  ferait chuter de partout & \sz{ne peut que} $\downarrow$ \\
2 & Topologie \x{} injectif & \textbf{M4} : un \code{config\_hash} \emph{rapporté}, fermant F2 &
  $\uparrow$ \sz{vers la 1\textsuperscript{re}} \\
3 & Liaison v4 \x{} ancre & \textbf{3b} : une ancre tierce -- \code{verify\_anchor} existe déjà &
  $\uparrow$ \sz{sur l'adversaire} \\
4 & Canari \x{} porte & retirer l'échappatoire, ou la rendre refusable par configuration &
  $\uparrow$ \sz{modeste} \\
5\,= & PII \x{} scellement & une policy tenant par \textbf{défaut fail-closed} & $\uparrow$ \sz{2 rangs} \\
5\,= & Plafonnement \x{} coût & rien de bon marché : sa faiblesse est la \emph{rareté} & \sz{stable} \\
\textbf{7} & \textbf{Co-signature R3} & \textbf{un producteur} servi + un registre peuplé &
  \best{$\uparrow$ six rangs} \\
8 & Revendication \x{} chute & fermer F1 : contraindre l'URL du SLM et décompter son egress &
  $\uparrow$ \sz{modeste} \\
\bottomrule
\end{tabular}
\end{center}

Un seul chantier déplace un couple de six rangs ; tous les autres valent un ou deux. Classés par
\emph{effet sur la valeur} plutôt que par difficulté ressentie, les travaux du portefeuille sont donc
écrasés par un seul item -- ce qui est exactement ce que mesurait l'écart de \pire{7}.

\section{5 · Les quatre documents de cette série, et lequel ouvrir}

Quatre analyses ont été produites en deux jours sur le même objet. Elles ne se remplacent pas -- elles
descendent -- et il vaut mieux dire laquelle répond à quoi que laisser un lecteur les ouvrir dans le
désordre.

\renewcommand{\arraystretch}{1.25}
\begin{center}\small
\begin{tabular}{@{}>{\raggedright\arraybackslash}p{5.0cm} >{\raggedright\arraybackslash}p{5.4cm} >{\raggedright\arraybackslash}p{5.6cm}@{}}
\toprule
\textbf{Le document} & \textbf{Il répond à} & \textbf{Ce qu'il faut en retenir} \\
\midrule
\code{Etat\_Des\_Lieux} \sz{(10 août)} &
  Où en est le dépôt ? &
  Tout est opt-in et \textbf{aucun profil ne l'allume} : un build par défaut ne sert aucun reçu \\
\code{Meilleurs\_Couples} \sz{(10 août)} &
  Quels couples existent ? &
  Les faits et les \textbf{faux couples}. Sa hiérarchie est périmée \sz{(lentille Kovant)}, ses faits
  ne le sont pas \\
\code{Couples\_Classement\_} \code{Intrinseque} \sz{(11 août)} &
  Dans quel ordre valent-ils, sans finalité ? &
  Le classement en vigueur, et pourquoi retirer l'axe extérieur a tout déplacé \\
\textbf{Ce document} \sz{(11 août)} &
  Comment marchent-ils exactement ? &
  Le mécanisme, le test qui échoue, et la \textbf{justification de chaque note} \\
\bottomrule
\end{tabular}
\end{center}

Deux remarques pour finir, qui portent sur la série plutôt que sur un document.

La première : \textbf{deux de ces quatre documents ont dû corriger le précédent}, et dans les deux cas
la correction est venue d'une commande écrite pour \emph{confirmer} une absence. C'est un taux d'erreur
élevé, et il n'est pas un accident de rédaction : les affirmations négatives -- « cela n'existe pas »,
« rien ne le produit » -- sont les seules qu'un lecteur ne peut pas vérifier en lisant, et donc les
seules qu'un auteur doit exécuter avant d'écrire. Les annexes de re-vérification existent pour cette
raison.

La seconde : \textbf{aucune cible d'usage n'est déclarée}, et aucun de ces documents n'en propose. Le
classement est sans finalité, ce qui est un état tenable -- c'était déjà la posture du 5 août. Si une
cible apparaît, la quatrième lentille se rejoue, et le §4 ci-dessus dit déjà de combien chaque rang
bougerait.

\begin{these}
\small\textbf{Portée et statut.} Mécanismes, emplacements, noms de tests et taxonomies sont des faits
lus dans le code du 11 août ; les \textbf{notes} sont un jugement structuré, pas une mesure. La chaîne
du reçu mesure \textbf{515 tests, 0 échec} \sz{(35 binaires)}, par exécution réelle. La couche nano et
\code{efe-gold-test-gate} n'ont été examinés dans aucun des quatre documents de cette série. Et le
rappel qui ne dépend d'aucune cible : le reçu atteste l'\textbf{intégrité} d'une suite de pas, jamais
la \textbf{correction} d'une décision, et il est \emph{tamper-evident pour un détenteur de clé}, pas
\emph{trustless pour un tiers}.
\end{these}

\end{document}
